\subsection{DIRECT DECOMPOSITION OF A RING}

Let A be a ring and \((b_{i})_{i\in I}\) a family of two-sided ideals of A. We shall call the homomorphism

\[
x \mapsto (\phi_ {i} (x)) _ {i \in I},
\]

where \(\phi_{i}\) is the canonical homomorphism of A onto \(\mathbf{A} / \mathfrak{b}_i\) , the canonical homomorphism of A into \(\prod_{i\in I}(\mathbf{A} / \mathfrak{b}_i)\) .

PROPOSITION 9. Let A be a ring and \((\mathfrak{b}_1, \ldots, \mathfrak{b}_n)\) two-sided ideals of A such that \(\mathfrak{b}_i + \mathfrak{b}_j = \mathbf{A}\) for \(i \neq j\) . The canonical homomorphism of A into \(\prod_{i=1}^{n} (\mathbf{A} / \mathfrak{b}_i)\) is surjective of kernel \(\bigcap_{i=1}^{n} \mathfrak{b}_i = \sum_{\sigma \in \mathfrak{E}_n} \mathfrak{b}_{\sigma(1)} \mathfrak{b}_{\sigma(2)} \ldots \mathfrak{b}_{\sigma(n)}\) .

Clearly the kernel is \(\bigcap_{i=1}^{n} \mathfrak{b}_i\) . To prove the surjectivity, it is necessary to show that, for every family \((x_i)_{1 \leqslant i \leqslant n}\) of elements of A, there exists \(x \in \mathbf{A}\) such that \(x \equiv x_i (\mathfrak{b}_i)\) for all \(1 \leqslant i \leqslant n\) . We prove this assertion by induction on the cardinal \(n\) of I, the case \(n \leqslant 1\) being trivial. By the induction hypothesis, there exists \(y \in \mathbf{A}\) such that \(y \equiv x_i (\mathfrak{b}_i)\) for \(1 \leqslant i \leqslant n-1\) . We look for an \(x\) of the

form \(y + z\) with \(z \in \mathbf{A}\) . Of necessity \(z \equiv 0(\mathfrak{b}_i)\) for \(i < n\) , that is \(z \in \mathfrak{b} = \bigcap_{i=1}^{\infty} \mathfrak{b}_i\) , and on the other hand \(z \equiv x_n - y(\mathfrak{b}_n)\) . Now \(\mathfrak{b}_n + \mathfrak{b} = A\) by no. 9, Proposition 6, whence the existence of \(z\) . Finally, the second expression of the kernel follows from no. 9, Proposition 7.

DEFINITION 7. Let A be a ring. A finite family \((\mathfrak{b}_i)_{i\in I}\) of two-sided ideals of A such that the canonical homomorphism of A into \(\prod_{i\in I}(\mathbf{A} / \mathfrak{b}_i)\) is an isomorphism is called a direct decomposition of A.

PROPOSITION 10. Let A be a ring, A' its centre and \((\mathfrak{b}_{i})_{i\in I}\) a finite family of two-sided ideals of A. The following conditions are equivalent:

\begin{enumerate}
\def\labelenumi{(\alph{enumi})}
\item
  the family \((\mathfrak{b}_i)_{i\in I}\) is a direct decomposition of A;
\item
  there exists a family \((e_i)_{i\in I}\) of idempotents of \(\mathbf{A}'\) such that \(e_i e_j = 0\) for \(i\neq j\) , \(1 = \sum_{i\in I}e_i\) and \(\mathfrak{b}_i = \mathbf{A}(1 - e_i)\) for \(i\in I\) ;
\item
  \(\mathfrak{b}_i + \mathfrak{b}_j = \mathbf{A}\) for \(i \neq j\) and \(\bigcap_{i \in I} \mathfrak{b}_i = \{0\}\) ;
\item
  \(\mathfrak{b}_i + \mathfrak{b}_j = \mathbf{A}\) for \(i \neq j\) and \(\prod_{i \in I} \mathfrak{b}_i = \{0\}\) for every total order on \(\mathbf{I}\) ;
\item
  there exists a direct decomposition \((\mathfrak{b}_i')_{i\in \mathbf{I}}\) of \(\mathbf{A}'\) such that \(\mathfrak{b}_i = \mathbf{Ab}_i'\) for \(i\in \mathbf{I}\) .
\end{enumerate}

(a) \(\Rightarrow\) (b). If condition (a) holds, A may be identified with the ring \(\prod_{i\in I}(\mathbf{A} / b_i)\)

and \(b_{i}\) with the kernel of \(pr_{i}\) . The existence of the \(e_{i}\) with the properties in (b) then follows from no. 10.

\begin{enumerate}
\def\labelenumi{(\alph{enumi})}
\setcounter{enumi}{1}
\tightlist
\item
  \(\Rightarrow\) (d). Suppose that the \(e_i\) exist with the properties in (b). For \(i \neq j\) , \(1 - e_i \in \mathfrak{b}_i\) , \(e_i = e_i(1 - e_j) \in \mathfrak{b}_j\) , hence \(1 \in \mathfrak{b}_i + \mathfrak{b}_j\) and \(\mathbf{A} = \mathfrak{b}_i + \mathfrak{b}_j\) . On the other hand, if I is given a total ordering and \((x_i)_{i \in I}\) is a family of elements of A, then, since the \(e_i\) are central,
\end{enumerate}

\[
\prod_ {i \in \mathrm{I}} x _ {i} (1 - e _ {i}) = \Bigl (\prod_ {i \in \mathrm{I}} x _ {i} \Bigr) \Bigl (\prod_ {i \in \mathrm{I}} (1 - e _ {i}) \Bigr) = \Bigl (\prod_ {i \in \mathrm{I}} x _ {i} \Bigr) \Bigl (1 - \prod_ {i \in \mathrm{I}} e _ {i} \Bigr) = 0
\]

hence \(\prod_{i\in I}b_i = \{0\}\) .

\begin{enumerate}
\def\labelenumi{(\alph{enumi})}
\setcounter{enumi}{3}
\item
  \(\Rightarrow\) (c). This follows from no. 9, Proposition 7.
\end{enumerate}

(c) \(\Rightarrow\) (a). This follows from Proposition 9.

Thus conditions (a), (b), (c) and (d) are equivalent. Suppose they hold. By \((\mathbf{b}) \Rightarrow (\mathbf{a})\) , the family of \(\mathfrak{b}_i' = \mathbf{A}'(1 - e_i)\) is a direct decomposition of \(\mathbf{A}'\) . Then \(\mathfrak{b}_i = \mathbf{A}(1 - e_i) = \mathbf{Ab}_i'\) for all \(i \in \mathbf{I}\) . Hence condition (e) holds.

Finally, suppose condition (e) holds. By (a) \(\Rightarrow\) (b), there exists a family \((e_i)_{i\in I}\) of idempotents of \(A'\) such that \(e_i e_j = 0\) for \(i \neq j\) , \(1 = \sum_{i\in I} e_i\) and \(\mathfrak{b}_i' = A'(1 - e_i)\) for \(i \in I\) . Then \(\mathfrak{b}_i = Abi' = A(1 - e_i)\) for \(i \in I\) , hence condition (b) holds.

Remark. Let A be a ring. Let \((a_{i})_{i\in I}\) be a finite family of subgroups of the additive group \(A^{+}\) of A such that \(A^{+}\) is the direct sum of the \(a_{i}\) . Suppose \(a_{i}a_{i}\subset ai\) for \(i\in I\) and \(a_{i}a_{j} = \{0\}\) for \(i\neq j\) . Then \(a_{i}\) is for all \(i\in I\) a two-sided ideal of A. With addition and multiplication induced by those on A, \(a_{i}\) is a ring with unit element the component of \(1\in A\) in \(a_{i}\) . If \(b_{i} = \sum_{j\neq i}a_{j}\) , clearly the \(b_{i}\) satisfy condition (c) of Proposition 10 and hence \((b_{i})_{i\in I}\) is a direct decomposition of A, which is said to be defined by \((a_{i})_{i\in I}\) .

\BourbakiRunInHeading{Example: Ideals and quotient rings of Z}

An ideal of \(\mathbf{Z}\) is an additive subgroup of \(\mathbf{Z}\) and hence of the form \(n\) . \(\mathbf{Z}\) with \(n \geqslant 0\) ; conversely, for every integer \(n \geqslant 0\) , the set \(n\) . \(\mathbf{Z}\) is an ideal, the principal ideal \((n)\) . Thus every ideal of \(\mathbf{Z}\) is principal and is represented uniquely in the form \(n\mathbf{Z}\) with \(n \geqslant 0\) . The ideal (1) is equal to \(\mathbf{Z}\) , the ideal (0) consists of 0 and the ideals distinct from \(\mathbf{Z}\) and \(\{0\}\) are therefore of the form \(n\mathbf{Z}\) with \(n > 1\) . If \(m \geqslant 1\) and \(n \geqslant 1\) , \(m\mathbf{Z} \supset n\mathbf{Z}\) if and only if \(n \in m\) . \(\mathbf{Z}\) , that is \(m\) divides \(n\) . Therefore, for the ideal \(n\mathbf{Z}\) to be maximal, it is necessary and sufficient that there exist no integer \(m > 1\) distinct from \(n\) and dividing \(n\) ; in other words, the maximal ideals of \(\mathbf{Z}\) are the ideals of the form \(p\mathbf{Z}\) where \(p\) is a prime number (§ 4, no. 10, Definition 16).

Let \(m\) and \(n\) be two integers \(\geqslant 1\) . The ideal \(m\mathbf{Z} + n\mathbf{Z}\) is principal, whence there is an integer \(d \geqslant 1\) characterized by \(d\mathbf{Z} = m\mathbf{Z} + n\mathbf{Z}\) ; for every integer \(r \geqslant 1\) , the relation ``r divides d'' is equivalent to \(rZ \supset dZ\) and hence to `` \(rZ \supset mZ\) and \(rZ \supset nZ\) '', that is to ``r divides m and n''. It is thus seen that the common divisors of m and n are the divisors of d and that d is the greatest of the divisors \(\geqslant 1\) common to m and n; d is called the greatest common divisor (abbreviated to g.c.d.) of m and n.~As \(dZ = mZ + nZ\) , there exist two integers x and y such that \(d = mx + ny\) . m and n are said to be relatively prime if their g.c.d. is equal to 1. It amounts to the same to assume that there exist integers x and y with \(mx + ny = 1\) .

The intersection of the ideals mZ and nZ is non-zero for it contains mn and hence is of the form rZ with \(r \geqslant 1\) . Arguing as above, it is seen that the multiples of r are the common multiples of m and n and that r is the least of the integers \(\geqslant 1\) which are common multiples of m and n; it is called the least common multiple (l.c.m.) of m and n.

The product of the ideals \(m\mathbf{Z}\) and \(n\mathbf{Z}\) is the set of \(\sum_{i=1}^{r} mx_i ny_i = mn\left(\sum_{i=1}^{r} x_i y_i\right)\) for \(x_1, \ldots, y_r \in \mathbf{Z}\) and hence is equal to \(mn\mathbf{Z}\) .

For every integer \(n \geqslant 1\) , the quotient ring Z/nZ is called the ring of integers modulo n; it has n elements, which are the classes modulo n of the integers 0, 1, 2, \(\ldots\) , n - 1. For n = 1, we obtain the zero ring.

PROPOSITION 11. Let \(n_1, \ldots, n_d\) be integers \(\geqslant 1\) which are relatively prime in pairs and \(n = n_1 \ldots n_r\) . The canonical homomorphism of \(\mathbf{Z}\) into the product ring \(\prod_{i=1}^{r} \mathbf{Z} / n_i \mathbf{Z}\) is surjective of kernel \(n\mathbf{Z}\) and defines a ring isomorphism of \(\mathbf{Z} / n\mathbf{Z}\) onto \(\prod_{i=1}^{r} \mathbf{Z} / n_i \mathbf{Z}\) .

Let \(a_{i} = n_{i}\mathbf{Z}\) for \(i = 1, \ldots, r\) . By hypothesis, \(a_{i} + a_{j} = \mathbf{Z}\) for \(i \neq j\) . The proposition then follows from Proposition 9.

The above results, as also those concerning decomposition into prime factors, will be generalized in Chapter VII, § 1, which is devoted to the study of principal ideal domains and in Commutative Algebra, Chapter VII, § 3, which is devoted to the study of factorial domains.
% semantic-fragments: legacy-input-seam
\relax{}
